\subsection{Analisis Kebutuhan Keamanan dan Data Provenance}

Ini menjawab RM3.

Referensi yang bisa digunakan:
\begin{itemize}
    \item Review data provenance di healthcare: pentingnya chain-of-custody, integritas, dan keterlacakan.
    
    \item Paper blockchain untuk PGHD / marketplace PGHD: blockchain digunakan untuk menjamin provenance dan auditable trail.
    
    \item Review blockchain untuk patient data security: transparansi, immutability, encryption, access control.
\end{itemize}

Hal yang kamu analisis:
\begin{itemize}
    \item Risiko jika PGHD tidak ada provenance:
    \begin{itemize}
        \item klinisi tidak percaya data,
        
        \item potensi manipulasi (user/spoofed device).
    \end{itemize}
    
    \item Keperluan:
    \begin{itemize}
        \item mengikat setiap entri PGHD ke identity kriptografis pasien,
        
        \item signature harus diverifikasi di smart contract IOTA (MoveVM),
        
        \item ledger hanya menyimpan hash/metadata, bukan data plaintext (privasi).
    \end{itemize}
\end{itemize}

Output analisis kebutuhan:
\begin{enumerate}
    \item N1: setiap entri PGHD memiliki patient id, timestamp, device/source metadata, dan signature,
    
    \item N2: smart contract harus bisa memverifikasi signature dan memastikan patient id terdaftar dalam DecMed,
    
    \item N3: ledger record harus immutable dan dapat diaudit.
\end{enumerate}